% ============================================================================
% CHAPTER 2: USER GUIDE - COMPREHENSIVE USAGE INSTRUCTIONS
% ============================================================================

\chapter{User Guide - End-to-End Workflows}

\section{Getting Started}

\subsection{Initial Login}

\begin{figure}[h]
\centering
\fbox{\textit{[SCREENSHOT PLACEHOLDER: Login Page with Username/Password Fields]}}
\caption{Enterprise Login Screen}
\label{fig:login}
\end{figure}

\textbf{Default Credentials:}
\begin{itemize}
    \item Username: \texttt{admin@example.com}
    \item Password: \texttt{admin123} (Change immediately!)
\end{itemize}

\begin{warningbox}
\textbf{Security Alert:} Default credentials must be changed on first login. Navigate to Settings → Security → Change Password.
\end{warningbox}

\subsection{Dashboard Overview}

\begin{figure}[h]
\centering
\fbox{\textit{[SCREENSHOT PLACEHOLDER: Main Dashboard with Key Metrics]}}
\caption{Dashboard - Home View}
\label{fig:dashboard}
\end{figure}

The dashboard provides real-time insights:
\begin{itemize}
    \item \textbf{Document Statistics} - Total, processed, pending
    \item \textbf{Classification Breakdown} - ICSR, Safety Info, Non-ICSR
    \item \textbf{Processing Queue} - Current workload
    \item \textbf{System Health} - Service status indicators
\end{itemize}

\section{Document Ingestion Workflow}

\subsection{Upload Local Documents (PDF)}

\textbf{Step-by-Step Process:}

\begin{enumerate}
    \item Navigate to \textbf{Project Spaces} from the sidebar
    \item Click \textbf{Create New Space} or select existing project
    \item Click \textbf{Upload Documents} button
    \item Configure ingestion options:
    \begin{itemize}
        \item Enable/Disable ICSR Classification
        \item Select Classification Model (Rule + BERT + LLM)
        \item Enable Entity Extraction
        \item Force Full Ensemble (optional)
    \end{itemize}
    \item Drag \& drop PDF files or click to browse
    \item Monitor upload progress
    \item View processed documents in Document Library
\end{enumerate}

\begin{figure}[h]
\centering
\fbox{\textit{[SCREENSHOT PLACEHOLDER: Project Spaces - Document Upload Interface]}}
\caption{Document Upload Interface}
\label{fig:upload}
\end{figure}

\begin{commandbox}
\textbf{Supported File Types:}
\begin{itemize}[noitemsep]
    \item PDF (up to 50 MB per file)
    \item Maximum 100 files per batch
    \item Text-based PDFs (OCR for scanned documents)
\end{itemize}
\end{commandbox}

\subsection{Ingest from PubMed}

\textbf{Configuration:}

\begin{lstlisting}[style=pythonstyle]
# Navigate to: Project Spaces → Configure Source → PubMed
source_config = {
    "source": "pubmed",
    "queries": [
        "COVID-19 adverse events",
        "vaccine safety pharmacovigilance"
    ],
    "max_results_per_query": 50,
    "date_range": {
        "start": "2024-01-01",
        "end": "2024-12-31"
    },
    "filters": {
        "article_types": ["Clinical Trial", "Case Reports"],
        "languages": ["eng"]
    }
}
\end{lstlisting}

\textbf{Execution Steps:}

\begin{enumerate}
    \item Click \textbf{Configure External Sources}
    \item Select \textbf{PubMed} from dropdown
    \item Enter search queries (one per line)
    \item Set maximum results per query (default: 50)
    \item Enable \textbf{ICSR Classification}
    \item Click \textbf{Start Ingestion}
    \item Monitor progress in real-time
\end{enumerate}

\begin{figure}[h]
\centering
\fbox{\textit{[SCREENSHOT PLACEHOLDER: PubMed Ingestion Configuration Panel]}}
\caption{PubMed Source Configuration}
\label{fig:pubmed-config}
\end{figure}

\textbf{Ingestion Status Tracking:}

\begin{table}[h]
\centering
\caption{Ingestion Status Indicators}
\begin{tabularx}{\textwidth}{lX}
\toprule
\textbf{Status} & \textbf{Description} \\
\midrule
\textcolor{warningyellow}{Pending} & Queued for processing \\
\textcolor{primaryblue}{Processing} & Currently being ingested \\
\textcolor{successgreen}{Completed} & Successfully processed \\
\textcolor{errorred}{Failed} & Error occurred (check logs) \\
\bottomrule
\end{tabularx}
\end{table}

\subsection{Batch Processing via Command Line}

For automated workflows or large-scale ingestion:

\begin{commandbox}
\begin{lstlisting}[style=bashstyle]
# Navigate to project directory
cd /path/to/data-ingestion-platform

# Run batch ingestion script
docker-compose exec api-gateway python scripts/batch_ingest.py \
    --source pubmed \
    --query "drug safety adverse events" \
    --max-results 500 \
    --enable-classification \
    --space-id <project_space_id>

# Monitor progress
docker-compose logs -f api-gateway | grep "Ingestion"
\end{lstlisting}
\end{commandbox}

\section{Document Classification}

\subsection{Classification System Overview}

The platform uses a \textbf{3-tier ensemble classification system}:

\begin{figure}[h]
\centering
\begin{tikzpicture}[
    node distance=2cm,
    startstop/.style={rectangle, rounded corners, minimum width=3cm, minimum height=1cm, text centered, draw=black, fill=primaryblue!30},
    process/.style={rectangle, minimum width=3cm, minimum height=1cm, text centered, draw=black, fill=successgreen!30},
    decision/.style={diamond, minimum width=3cm, minimum height=1cm, text centered, draw=black, fill=warningyellow!30},
]
    \node (start) [startstop] {Document Input};
    \node (tier1) [process, below of=start] {Tier 1: Rule-Based};
    \node (check1) [decision, below of=tier1] {Confidence > 0.85?};
    \node (tier2) [process, right of=check1, xshift=3cm] {Tier 2: BERT};
    \node (check2) [decision, below of=tier2] {Confidence > 0.75?};
    \node (tier3) [process, below of=check2] {Tier 3: LLM (GPT-4)};
    \node (ensemble) [process, left of=tier3, xshift=-3cm] {Ensemble Voting};
    \node (output) [startstop, below of=ensemble] {Final Classification};
    
    \draw [->] (start) -- (tier1);
    \draw [->] (tier1) -- (check1);
    \draw [->] (check1) -- node[anchor=south] {No} (tier2);
    \draw [->] (check1) -- node[anchor=east] {Yes} (ensemble);
    \draw [->] (tier2) -- (check2);
    \draw [->] (check2) -- node[anchor=east] {No} (tier3);
    \draw [->] (check2) -- node[anchor=south] {Yes} (ensemble);
    \draw [->] (tier3) -- (ensemble);
    \draw [->] (ensemble) -- (output);
\end{tikzpicture}
\caption{3-Tier Ensemble Classification Flow}
\label{fig:classification-flow}
\end{figure}

\subsection{Classification Categories}

\begin{table}[h]
\centering
\caption{Classification Categories \& Definitions}
\begin{tabularx}{\textwidth}{lX}
\toprule
\textbf{Category} & \textbf{Definition \& Criteria} \\
\midrule
\textbf{ICSR} & Individual Case Safety Report - contains specific patient case with adverse event, drug, and temporal relationship \\
\textbf{Safety\_Information} & General safety information, aggregate data, clinical trial results without individual patient details \\
\textbf{Non-ICSR} & Medical literature without safety relevance (e.g., efficacy studies, mechanism of action) \\
\textbf{Rejected} & Not relevant to pharmacovigilance (e.g., financial reports, opinion pieces) \\
\bottomrule
\end{tabularx}
\end{table}

\subsection{Viewing Classification Results}

\textbf{Document Library View:}

\begin{figure}[h]
\centering
\fbox{\textit{[SCREENSHOT PLACEHOLDER: Document Library with Classification Badges]}}
\caption{Document Library - Classification Results}
\label{fig:doc-library}
\end{figure}

Each document displays:
\begin{itemize}
    \item \textbf{Classification Badge} - Color-coded by category
    \item \textbf{Confidence Score} - 0-100\% reliability
    \item \textbf{Model Attribution} - Which model made the decision
    \item \textbf{Processing Time} - Milliseconds taken
    \item \textbf{Cost} - API cost (for LLM tier)
\end{itemize}

\subsection{Force Full Ensemble Mode}

To run all three models for comparison or validation:

\begin{enumerate}
    \item In Project Spaces, enable \textbf{"Force Full Ensemble"}
    \item This bypasses confidence thresholds
    \item All models (Rule + BERT + LLM) run in parallel
    \item Final result determined by weighted voting
    \item Useful for: Quality assurance, model comparison, edge cases
\end{enumerate}

\begin{infobox}
\textbf{Cost Consideration:} Force Full Ensemble mode increases API costs due to LLM invocation for every document. Use judiciously for critical documents only.
\end{infobox}

\section{Knowledge Hub - 3D Graph Visualization}

\subsection{Accessing the Knowledge Hub}

Navigate to \textbf{Knowledge Hub} from the main sidebar.

\begin{figure}[h]
\centering
\fbox{\textit{[SCREENSHOT PLACEHOLDER: Knowledge Hub - 3D WebGL Graph View]}}
\caption{Knowledge Hub - Interactive 3D Graph}
\label{fig:knowledge-hub}
\end{figure}

\subsection{Graph Rendering Modes}

\textbf{Toggle Between Modes:}

\begin{table}[h]
\centering
\caption{Graph Rendering Modes Comparison}
\begin{tabularx}{\textwidth}{lXX}
\toprule
\textbf{Mode} & \textbf{2D Canvas} & \textbf{3D WebGL} \\
\midrule
Rendering & HTML5 Canvas & Three.js WebGL \\
Performance & 30-45 FPS & 60 FPS \\
Max Nodes & 500 & 5000+ \\
Interactions & Drag, Zoom & Drag, Zoom, Rotate \\
Lighting & None & PBR Lighting \\
Best For & Quick overview & Deep exploration \\
\bottomrule
\end{tabularx}
\end{table}

\subsection{3D WebGL Controls}

\textbf{Mouse/Trackpad Controls:}
\begin{itemize}
    \item \textbf{Left-Click + Drag} - Rotate graph in 3D space
    \item \textbf{Scroll Wheel} - Zoom in/out
    \item \textbf{Right-Click + Drag} - Pan camera
    \item \textbf{Click Node} - Select and highlight connections
    \item \textbf{Double-Click Node} - Center and zoom to node
\end{itemize}

\textbf{Button Controls:}
\begin{itemize}
    \item \faSearchPlus\ \textbf{Zoom In} - Reduce camera distance by 30\%
    \item \faSearchMinus\ \textbf{Zoom Out} - Increase camera distance by 30\%
    \item \faBullseye\ \textbf{Center \& Fit} - Auto-fit all nodes with animation
    \item \faRedo\ \textbf{Restart Layout} - Reheat physics simulation
    \item \faExpand\ \textbf{Fullscreen} - Toggle fullscreen mode
\end{itemize}

\begin{figure}[h]
\centering
\fbox{\textit{[SCREENSHOT PLACEHOLDER: 3D Graph Controls Panel (Right Side)]}}
\caption{Professional Control Panel}
\label{fig:graph-controls}
\end{figure}

\subsection{Advanced Filtering}

\textbf{Progressive Disclosure:}
\begin{itemize}
    \item Default: Show top 50 most connected nodes
    \item Toggle \textbf{"Show All Nodes"} to display full graph
    \item Filter by minimum connections (0-10)
    \item Filter by confidence threshold (0-100\%)
\end{itemize}

\textbf{Community Detection:}
\begin{itemize}
    \item Enable \textbf{"Show Communities"} toggle
    \item Nodes grouped by Louvain algorithm
    \item Color-coded by semantic clusters
    \item Useful for identifying drug-disease relationships
\end{itemize}

\begin{figure}[h]
\centering
\fbox{\textit{[SCREENSHOT PLACEHOLDER: Knowledge Hub with Community Detection Enabled]}}
\caption{Community Detection - Semantic Clustering}
\label{fig:communities}
\end{figure}

\subsection{Entity Details Panel}

Clicking any node opens a detailed drawer:

\begin{itemize}
    \item \textbf{Entity Information}
    \begin{itemize}
        \item Name, Type (Drug/Disease/Person/Org)
        \item Confidence score
        \item Source documents count
    \end{itemize}
    \item \textbf{Source Documents} (up to 20)
    \begin{itemize}
        \item Document title and snippet
        \item Highlighted entity mentions
        \item Click to open full document
    \end{itemize}
    \item \textbf{Interactive Chat}
    \begin{itemize}
        \item Ask questions about the entity
        \item AI-powered responses with citations
        \item Example: "What are the side effects?"
    \end{itemize}
    \item \textbf{Pin to Investigation Board}
    \begin{itemize}
        \item Save entity for collaborative review
        \item Add to active safety signal investigations
    \end{itemize}
\end{itemize}

\begin{figure}[h]
\centering
\fbox{\textit{[SCREENSHOT PLACEHOLDER: Entity Details Drawer - Source Documents \& Chat]}}
\caption{Entity Details Panel}
\label{fig:entity-drawer}
\end{figure}

\section{ICSR Reviewer Workflow}

\subsection{Automated ICSR Extraction}

For documents classified as ICSR, the system automatically extracts:

\begin{table}[h]
\centering
\caption{ICSR Extracted Fields}
\begin{tabularx}{\textwidth}{lX}
\toprule
\textbf{Category} & \textbf{Extracted Fields} \\
\midrule
Patient & Age, Gender, Weight, Initials \\
Drug & Product Name, Dosage, Route, Start/End Date \\
Event & Adverse Event Description, Onset Date, Outcome, Seriousness \\
Reporter & Reporter Type, Country, Contact Info \\
Narrative & AI-generated summary, Timeline \\
\bottomrule
\end{tabularx}
\end{table}

\subsection{Review Interface}

\begin{figure}[h]
\centering
\fbox{\textit{[SCREENSHOT PLACEHOLDER: ICSR Reviewer - Side-by-Side View]}}
\caption{ICSR Reviewer - Document \& Extracted Data}
\label{fig:icsr-reviewer}
\end{figure}

\textbf{Review Process:}
\begin{enumerate}
    \item Navigate to \textbf{ICSR Reviewer}
    \item Select document from queue
    \item Review extracted fields in form
    \item Edit/correct fields as needed
    \item Add reviewer notes
    \item Mark as: Approved / Needs Revision / Rejected
    \item Submit for regulatory reporting
\end{enumerate}

\textbf{Validation Rules:}
\begin{itemize}
    \item Required fields highlighted in red
    \item Date format validation (ISO 8601)
    \item Drug name verification against database
    \item Seriousness criteria check
\end{itemize}

\section{Investigation Boards}

\subsection{Creating an Investigation Board}

\begin{enumerate}
    \item Navigate to \textbf{Investigation Boards}
    \item Click \textbf{Create New Board}
    \item Enter board details:
    \begin{itemize}
        \item Title: "Cardiac Events - Drug X"
        \item Description: Objective and scope
        \item Priority: Low / Medium / High / Critical
        \item Due Date: Target completion date
    \end{itemize}
    \item Assign team members
    \item Click \textbf{Create Board}
\end{enumerate}

\begin{figure}[h]
\centering
\fbox{\textit{[SCREENSHOT PLACEHOLDER: Investigation Board Creation Modal]}}
\caption{Create New Investigation Board}
\label{fig:create-board}
\end{figure}

\subsection{Pinning Entities \& Documents}

\textbf{From Knowledge Hub:}
\begin{itemize}
    \item Click any entity node
    \item Click \textbf{Pin to Investigation Board}
    \item Select target board
    \item Add pin reason/notes
\end{itemize}

\textbf{From Document Viewer:}
\begin{itemize}
    \item Open any document
    \item Click \textbf{Pin to Board} button
    \item Select entities to include
    \item Choose investigation board
\end{itemize}

\subsection{Collaborative Features}

\begin{itemize}
    \item \textbf{Real-time Updates} - Changes visible to all team members
    \item \textbf{Comments} - Add notes on any pinned item
    \item \textbf{Task Assignment} - Assign review tasks to users
    \item \textbf{Status Tracking} - Mark items as: To Do / In Progress / Completed
    \item \textbf{Export} - Generate PDF reports with all evidence
\end{itemize}

\begin{figure}[h]
\centering
\fbox{\textit{[SCREENSHOT PLACEHOLDER: Investigation Board - Collaborative Workspace]}}
\caption{Investigation Board - Active Case}
\label{fig:investigation-board}
\end{figure}

\section{Monitoring \& Performance}

\subsection{System Dashboard (Grafana)}

Access: \texttt{http://localhost:3001}

\begin{figure}[h]
\centering
\fbox{\textit{[SCREENSHOT PLACEHOLDER: Grafana Dashboard - System Metrics]}}
\caption{Grafana - Real-time Performance Monitoring}
\label{fig:grafana}
\end{figure}

\textbf{Key Metrics Monitored:}
\begin{itemize}
    \item CPU \& Memory usage per service
    \item Request rate \& response times
    \item Database query performance
    \item Classification throughput (docs/hour)
    \item Error rates \& exception tracking
    \item Queue depth \& processing lag
\end{itemize}

\subsection{Audit Logs}

Navigate to \textbf{Settings → Security → Audit Logs}

\begin{itemize}
    \item All user actions logged with timestamps
    \item Searchable by: User, Action, Date Range
    \item Exportable to CSV for compliance
    \item Retention policy: 90 days default
\end{itemize}

\section{User Management \& Security}

\subsection{Role-Based Access Control (RBAC)}

\begin{table}[h]
\centering
\caption{User Roles \& Permissions}
\begin{tabularx}{\textwidth}{lX}
\toprule
\textbf{Role} & \textbf{Permissions} \\
\midrule
\textbf{Admin} & Full system access, user management, security settings \\
\textbf{Data Manager} & Upload documents, configure sources, manage projects \\
\textbf{Reviewer} & Review ICSRs, approve classifications, add annotations \\
\textbf{Analyst} & View documents, search, knowledge hub (read-only) \\
\textbf{Guest} & Limited access, cannot modify data \\
\bottomrule
\end{tabularx}
\end{table}

\subsection{Adding Users}

\begin{enumerate}
    \item Navigate to \textbf{Settings → User Management}
    \item Click \textbf{Add New User}
    \item Fill in user details:
    \begin{itemize}
        \item Full Name
        \item Email Address
        \item Assign Role
        \item Set Password (or send invite link)
    \end{itemize}
    \item Click \textbf{Create User}
    \item User receives email invitation
\end{enumerate}

\begin{figure}[h]
\centering
\fbox{\textit{[SCREENSHOT PLACEHOLDER: User Management Interface]}}
\caption{User Management Dashboard}
\label{fig:user-mgmt}
\end{figure}

\subsection{Multi-Factor Authentication (MFA)}

\textbf{Enable MFA for Enhanced Security:}
\begin{enumerate}
    \item Navigate to \textbf{Settings → Security → MFA}
    \item Click \textbf{Enable MFA}
    \item Scan QR code with authenticator app (Google Authenticator, Authy)
    \item Enter verification code
    \item Save backup codes in secure location
\end{enumerate}

\begin{warningbox}
\textbf{Backup Codes:} Store backup codes securely. If authenticator is lost, these codes are the only way to regain access.
\end{warningbox}

\end{document}

